\scanpage{8}

\mspage{4)}
Cette composition est associative à iso can. près, et elle a des identités bilatères (\deleted{\unsure{pour|par}} \unsure{un tel|cet|unique|un} iso. can. !) savoir

\diagram{\begin{tikzcd}[ampersand replacement=\&,column sep=small,row sep=small]
x \arrow[dr,"1"'] \& x \arrow[d,tail,"1","\sim"'] \\
{} \& x
\end{tikzcd}}

noté $\underline{1}_x$

On a d'autre part des \deleted{\unsure{applic}} foncteurs

\[
 \operatorname{Hom}_{C}(x,y)_{\rm discret}\longrightarrow\underline{\mathrm{Ch}}_0(x,y)
\]
\deleted{compatibles avec}
\[
 f\longmapsto (f,\mathrm{id}_x)
\]

compatibles aux compositions. Ceci posé, soit

\[
 \pi_0(x,y)\overset{\rm déf}{=}\pi_0\bigl(\underline{\mathrm{Ch}}_0(x,y)\bigr),
\]

alors la composit. des $\underline{\mathrm{Ch}}_0$ définit une composition des $\pi_0(x,y)$, qui est associative (str\textsuperscript{t} !) et admet les classes des $\underline{1}_a$ comme unités bilatères. D'où une catégorie, que je note $\widetilde C^{\unsure{2|z|?}}$, et un foncteur

\[
 C\longrightarrow\widetilde C
\]

Pour l'instant, on n'a utilisé rien sur $(C,W)$, [\deleted{sauf} que \unsure{$\subset$|la} la donnée des $\mathrm{cof}_{-}\subset\mathrm{Fl}\,C$, on convient de prendre les flèches $i$ dans \deleted{\unsure{cof triv|cofibs}} \inserted{$\mathrm{couniv.}W$} -- mais \deleted{\unsure{on}} on peut prendre la construction, une cof \inserted{pas néc. $=\mathrm{cof}_{W}$,} telle que $\mathrm{cof}\cap W\subset\deleted{\unsure{W}}\mathrm{Couniv.}\,W$\,]. \deleted{Mais,} On veut prouver que ce foncteur transforme $f\in W$ en iso.

\diagram{\begin{tikzcd}[ampersand replacement=\&,column sep=small,row sep=small]
x \arrow[r,"f","\sim"'] \& y \arrow[dr,"{g?}"'] \& x \arrow[d,tail,"{i?}","\sim"'] \\
{} \& {} \& \bar x
\end{tikzcd}}

on cherche $g,i$ tels que $(gf,i)\in\mathrm{comp.}\ \mathrm{conn.}$ de $\underline{1}_x$ dans $\pi_0(\underline{\mathrm{Ch}}_0(x,x))$ et $(\bar f g,j)\in\mathrm{comp.}\ \mathrm{conn.}$ de $\underline{1}_y$

\diagram{\begin{tikzcd}[ampersand replacement=\&,column sep=small,row sep=small]
y \arrow[dr,"g"'] \& x \arrow[r,"f","\sim"'] \arrow[d,tail,"i","\sim"'] \& y \arrow[d,tail,"j"'," \sim"] \\
{} \& \bar x \arrow[r,"{\text{(cocart)}\ \sim}","{\bar f\ \sim}"'] \& \bar y
\end{tikzcd}}

% archive sheet number: 7
